\begin{titlepage}
\thispagestyle{empty}
\centering
\vspace*{1.0in}
{\small\color{teal}\textls[160]{\textsc{A compact field guide}}\par}
\vspace{0.45in}
{\fontsize{34}{39}\selectfont\color{navy}\textsc{Book\\Title}\par}
\vspace{0.28in}
{\Large\itshape One-line subtitle\par}
\vspace{0.5in}
{\color{gold}\rule{0.55\linewidth}{1.2pt}\par}
\vspace{0.45in}
{\large Context line: course, goal, or reader\par}
\vfill
{\small Version 1.0 \enspace--\enspace MONTH YEAR\par}
\end{titlepage}

\begin{epubonly}
\chapter*{About this book}
\end{epubonly}

\thispagestyle{empty}
\vspace*{\fill}
\noindent\small
This is a deliberately short learning text, not an encyclopaedia. The core chapters
assume PREREQUISITES; Appendix A provides an intuition-first, self-contained refresher.
The text states the conditions behind its central identities and points to primary
references when a complete proof would defeat the purpose of a compact guide.

\medskip
\noindent The source is designed to be revised. If you find an error, record the smallest
counterexample you can.
\vspace*{\fill}
\clearpage

\tableofcontents
\clearpage

\chapter*{Reader's contract}
\addcontentsline{toc}{chapter}{Reader's contract}

State the handful of questions that persist across the whole subject, and how to read
the book (two passes; do the checkpoints; answers at the back).

\section*{Map to the course}

\begin{center}
\small
\begin{tabularx}{\textwidth}{@{}lXX@{}}
\toprule
\textbf{Read} & \textbf{Lecture} & \textbf{Lasting use} \\
\midrule
Ch. 1 & 1--2 & Problem formulation \\
Ch. 2 & 3 & The reusable core \\
App. A & as needed & Refresher \\
\bottomrule
\end{tabularx}
\end{center}

\begin{takeaway}
The book is complete when you can reconstruct its spine, not when you remember every
equation.
\end{takeaway}
